\section{March 03, 2026:...}

\paragraph{Plan}
\begin{itemize}
    \item [I] Basis of algebraic cycles:
    algebraic cycles, rational equivalence, Chow groups, chern and Segre classes, Gysin maps, intersection products, Grothendieck-Riemann-Roch, Motive.(Fulton)
    \item[II] Higher Chow groups and motivic cohomology (Bloch, Mazza-Voevodsky-Weibel)
\end{itemize}

\subsection{Rational equivalence}

    A scheme will mean a separate scheme of finite type over a field \(\kk\) unless otherwise specified. 
    A variety will mean an integral scheme.

    \begin{example}
        The closed subscheme \(X\) of \(\bbA^n\) are in one-to-one correspondence with ideals \(I\) of \(\kk[x_1,\dots,x_n]\).
    \end{example}

    \begin{example}
        A base change of a variety \(X\) may not be a variety. 
        For example, let \(X = \Spec \bbR[X]/(X^2+1)\) be a variety over \(\bbR\). 
        Then the base change \(X_{\bbC} = X \times_{\Spec \bbR} \Spec \bbC\) is not a variety since it is not irreducible.
    \end{example}

    \begin{definition}
        An \(i\)-cycle on a scheme \(X\) is a formal finite sum \(\sum n_j [V_j]\) where \(n_j \in \bbZ\) and \(V_j\) are closed subvarieties of \(X\) of dimension \(i\).
    \end{definition}

    \begin{definition}
        \[ Z_i(X) \coloneqq \{\text{all } i\text{-cycles on } X\} \]
        is the free abelian group generated by the \(i\)-dimensional closed subvarieties of \(X\).
    \end{definition}

    \begin{definition}
        Let \(X\) be a variety of dimension \(n\) and \(f \in \fkk(X)^{\times}\) be a nonzero rational function on \(X\). 
        Then define the divisor of \(f\) to be
        \[ \divisor(f) \coloneqq \sum_{V} \ord_V(f) [V] \in Z_{n-1}(X) \]
        where the sum is taken over all codimension 1 subvarieties \(V\) of \(X\) and \(\ord_V(f)\) is the order of vanishing of \(f\) along \(V\) after normalizing \(X\).
        Explicitly, if \(\nu: \tilde{X} \to X\) is the normalization of \(X\), then 
        \[ \ord_V(f) \coloneqq \sum_{\tilde{V} \subset \tilde{X}, \nu(\tilde{V}) = V} [\fkk(\tilde{V}):\fkk(V)] \ord_{\tilde{V}}(f) \]
        where the sum is taken over all codimension 1 subvarieties \(\tilde{V}\) of \(\tilde{X}\) that map to \(V\) under \(\nu\).
    \end{definition}

    \begin{example}
        Let \(f \in \kk[t]\) be a nonzero polynomial of degree \(d\). 
        By prime factorization, write \(f = \prod_{i=1}^r f_j^{m_j}\) where \(f_j\) are distinct irreducible polynomials and \(m_j\) are positive integers. 
        Regard \(f\) as a rational function on \(\bbP^1\)。
        Then the divisor of \(f\) is
        \[ \divisor(f) = \sum_{j=1}^r m_j [V(f_j)] - d [\infty], \]
        where \(V(f_j)\) is the closed point of \(\bbP^1\) defined by \(f_j\) and \(\infty\) is the point at infinity of \(\bbP^1\).
    \end{example}

    \begin{definition}
        Let \(X\) be a scheme. 
        A \(i\)-cycles \(\alpha \in Z_i(X)\) are \emph{rationally equivalent to zero} 
        if there exists finitely many subvarieties \(W_j\) of \(X\) of dimension \(i+1\) and rational functions \(f_j \in \fkk(W_j)^{\times}\) such that
        \[ \alpha = \sum_j \divisor(f_j). \]
    \end{definition}

    For example, consider the \(0\)-cycle on \(X\).
    There are finitely many closed points \(P_j\) of \(X\) and integers \(n_j\) such that \(\alpha = \sum n_j [P_j]\). 
    Then \(\alpha\) is rationally equivalent to zero if and only if there are finitely many curves \(C_k\) on \(X\) and rational functions \(f_k \in \fkk(C_k)^{\times}\) such that
    \[ \alpha = \sum_k \divisor(f_k). \]

    \begin{definition}
        The \emph{Chow group} of \(i\)-cycles on \(X\) is defined to be the quotient group
        \[ \CH_i(X) \coloneqq Z_i(X) / \{\text{rationally equivalent to zero}\}. \]
        When \(X\) is smooth, we also write \(\CH^i(X) \coloneqq \CH_{n-i}(X)\) where \(n = \dim X\).
    \end{definition}

    \begin{example}
        \[ \CH_{n-1}(X) \cong \Cl(X). \]
        \[ \CH_{*}(X) \cong \CH_*(X_{\text{red}}). \]
    \end{example}

    We will study the covariance and contravariance of Chow groups.

    \paragraph{Proper pushforwards}
    Let \(f: X \to Y\) be a proper morphism of schemes, \(V \subset X\) a subvariety of dimension \(i\), and \(W = f(V)\) the image of \(V\) in \(Y\). 
    Define 
    \[ \deg(V/W) \coloneqq \begin{cases}
        [\fkk(V):\fkk(W)], & \text{if } \dim V = \dim W, \\
        0, & \text{otherwise}.
    \end{cases} \]
    For example, when \(\kk=\kkk\) of characteristic zero, if \(\dim W = \dim V\), then \(\deg(V/W)\) is the number of points in the fiber of \(f|_V: V \to W\) over a general point of \(W\).
    Define the pushforward of \([V]\) to be
    \[ f_*[V] \coloneqq \deg(V/W) [W]. \]
    Extend this definition linearly to get a homomorphism \(f_*: Z_i(X) \to Z_i(Y)\). 
    If \(g: Y \to Z\) is another proper morphism, then \(g_* \circ f_* = (g \circ f)_*\) because \(\deg(V/(g \circ f)(V)) = \deg(V/W) \deg(W/g(W))\).

    \begin{example}
        Let \(X\) be a scheme and \(f:X \to \Spec \kk\) be the structure morphism. 
        Define the degree map
        \[ \deg: Z_0(X) \to Z_0(\Spec \kk) \cong \bbZ, \quad \sum n_j [P_j] \mapsto \sum n_j [\rkk(P_j):\kk]. \]
    \end{example}

    \begin{theorem}
        Let \(f: X \to Y\) be a proper morphism and \(\alpha \in Z_i(X)\).
        If \(\alpha\) is rationally equivalent to zero on \(X\), then \(f_*\alpha\) is rationally equivalent to zero on \(Y\).
        In particular, \(f_*\) induces a homomorphism \(f_*: \CH_i(X) \to \CH_i(Y)\), i.e., \(\CH_*\) is covariant w.r.t. proper morphisms.
    \end{theorem}
    \begin{proof}
        May assume \(\alpha = \divisor(r)\) for some rational function \(r\) on a subvariety \(V\) of \(X\) of dimension \(i+1\).
        WTS: \(f_*\alpha \sim 0\).
        Replace \(X\) by \(V\) and \(Y\) by \(W \coloneqq f(V)\). 
        Then \(f: V \to W\) is a proper surjective morphism of varieties.
        The theorem follows from \cref{prop:pushforward_of_divisor_of_rational_funcitons}.
    \end{proof}

    \begin{proposition}\label{prop:pushforward_of_divisor_of_rational_funcitons}
        Let \(f: X \to Y\) be a surjective proper morphism of varieties and \(r \in \fkk(X)^{\times}\) a rational function. 
        Then
        \begin{enumerate}
            \item \(f_*\divisor(r) = 0\) if \(\dim W < \dim V\);
            \item \(f_*\divisor(r) = \divisor(\Nm_{X/Y}(r))\) if \(\dim W = \dim V\).
        \end{enumerate}
    \end{proposition}
    \begin{proof}[Proof of (b)]
        Let 
        \[ S \coloneqq \{y \in W: \dim f^{-1}(y) > 0\} \subset Y.\]
        Then \(\codim_Y S \geq 2\) since if \(\codim_Y S = 1\), then \(f^{-1}(S) = X\) which contradicts the surjectivity of \(f\).
        Replace \(Y\) by \(Y \setminus S\) and \(X\) by \(f^{-1}(Y \setminus S)\). 
        We may assume \(f: X \to Y\) is a finite morphism of varieties.

        Let \(g:\tilde{X} \to X\) and \(h:\tilde{Y} \to Y\) be the normalizations of \(X\) and \(Y\) respectively. 
        Then we have the induced morphism \(\tilde{f}: \tilde{X} \to \tilde{Y}\) such that \(h \circ \tilde{f} = f \circ g\).
        Then \(f, g, h, \tilde{f}\) are all finite morphisms of varieties and \(g,h\) are birational morphisms.
        We show that 
        \begin{itemize}
            \item [(b1)] for any rational function \(r\) on \(X\), we have \(g_*\divisor(r) = \divisor(r)\);
            \item [(b2)] \(\tilde{f}_*\divisor(r) = \divisor(\Nm_{X/Y}(r))\).
        \end{itemize}
        If (b1) and (b2) hold, then
        \[ f_*\divisor(r) = f_* g_* \divisor(r) = h_* \tilde{f}_* \divisor(r) = h_* \divisor(\Nm_{X/Y}(r)) = \divisor(\Nm_{X/Y}(r)). \]
        Thus we are done.

        For (b1), by definition, for a codimension 1 subvariety \(V\) of \(X\), we have
        \[ \ord_V(r) = \sum_{\tilde{V} \subset \tilde{X}, g(\tilde{V}) = V} [\fkk(\tilde{V}):\fkk(V)] \ord_{\tilde{V}}(r). \]
        Taking the sum over all codimension 1 subvarieties \(V\) of \(X\), we get
        \[ \divisor(r) = g_*\divisor(r). \]

        For (b2), let \(V\) be a codimension 1 subvariety of \(\tilde{Y}\). 
        Then we need to show that
        \begin{equation}\label{eq:the_star}
            \ord_V(\Nm_{X/Y}(r)) = \sum_{\tilde{V} \subset \tilde{X}, \tilde{f}(\tilde{V}) = V} [\fkk(\tilde{V}):\fkk(V)] \ord_{\tilde{V}}(r).
        \end{equation} 
        Set \(A = \calO_{\tilde{Y},V}\) and \(B\) be the integral closure of \(A\) in \(\fkk(X)\).
        Then \(A\) is a DVR and \(B\) is f.g. torsion-free \(A\)-module.
        Hence \(B \cong A^d\) for \(d = [\fkk(X):\fkk(Y)]\).
        Note that the maximal ideals of \(B\) are in one-to-one correspondence with the codimension 1 subvarieties \(\tilde{V}\) of \(\tilde{X}\) that map to \(V\) under \(\tilde{f}\). 
        May assume \(r \in B\), then \(B/rB \cong \prod_{\frakq \subset B} B_\frakq / r B_\frakq\) is an artinian ring, where the product is taken over all maximal ideals \(\frakq\) of \(B\).
        The RHS of \eqref{eq:the_star} is equal to \(\dim_{\fkk(V)} B/rB\).
        By the theory of elementary divisors, 
        \[ \begin{tikzcd}
            B \ar[r, "\text{times by } r"] \ar[d, "\cong"] & B \ar[d, "\cong"] \\
            A^d \ar[r, "M"] & A^d,
        \end{tikzcd} \]
        where \(M\) is a diagonal matrix.
        Then 
        \[ \dim_{\fkk(V)} B/rB = \dim_{\fkk(V)} A^d / (\det M) A^d = \text{ LHS of \eqref{eq:the_star}}. \]
        We are done.
    \end{proof}
    \begin{proof}[Proof of (a)]
        If \(\dim Y < \dim X - 1\), the assertion is obvious.
        Assume \(\dim Y = \dim X - 1\).
        Replacing \(Y\) by \(\Spec \fkk(Y) = \Spec \KK\) and \(X\) by \(X \times_Y \Spec \fkk(Y)\), we may assume \(Y = \Spec \fkk(Y)\) is a point and \(X\) is a proper curve.
        WTS: \(f_*\divisor(r) = 0\), i.e., \# of zeros of \(r\) equals \# of poles of \(r\) counting multiplicity.
        
        Let \(g: \tilde{X} \to X\) be the normalization of \(X\) and \(\tilde{f}: \tilde{X} \to Y\) be the induced morphism.
        Since \(g_*\divisor(r) = \divisor(r)\), we only need to show that \(\tilde{f}_*\divisor(r) = 0\).
        Choose any non-constant rational function \(t\) on \(\tilde{X}\).
        It gives a finite morphism \(t: \tilde{X} \to \bbP^1\).
        By (b), \(\tilde{f}_*\divisor(r) = p_*\divisor(\Nm_{X/Y}(r))\), where \(p: \bbP^1 \to \Spec \KK\) is the structure morphism. 

        Write \(\divisor(\Nm_{X/Y}(r)) = a/b\) with \(a,b\) coprime polynomials in \(\KK[t]\). 
        By prime factorization, \(a = \prod a_j^{m_j}\), \(\div(a) = \sum m_j [V(a_j)] - (\sum m_j[\KK(V(a_j)):\KK]) [\infty]\).
        This means 
        \[p_*\divisor(a) = \sum m_j [\KK(V(a_j)):\KK] - (\sum m_j[\KK(\infty):\KK]) [\KK(\infty):\KK] = 0. \]
        Similarly, \(p_*\divisor(b) = 0\).
        Hence \(p_*\divisor(\Nm_{X/Y}(r)) = p_*\divisor(a) - p_*\divisor(b) = 0\).
    \end{proof}

    \begin{definition}
        Let \(X\) be a proper scheme over \(\kk\), \(f:X \to \Spec \kk\) be the structure morphism.
        Define the degree map
        \[ \deg: \CH_0(X) \to \CH_0(\Spec \kk) \cong \bbZ, \quad \sum n_j [P_j] \mapsto \sum n_j [\rkk(P_j):\kk]. \]
    \end{definition}

    \begin{example}
        Properness is necessary for the degree map of \(\CH_0\) to be well-defined.
        Let \(X = \bbA^1\), \(t \in \fkk(\bbA^1) = \kk(t)\) be the coordinate function.
        Let \(P = V(t) \subset \bbA^1\) be the divisor of \(t\). 
        On the one hand, \(\deg([P]) = [\rkk(P):\kk] = 1\).
        On the other hand, \(P = \divisor(t)\) is rationally equivalent to zero, so \(\deg([P]) = 0\) if the degree map is well-defined.
    \end{example}
    